
\documentclass{article}
\usepackage{colortbl}
\usepackage{makecell}
\usepackage{multirow}
\usepackage{supertabular}

\begin{document}

\newcounter{utterance}

\centering \large Interaction Transcript for game `codenames', experiment `ambiguous', episode 1 with reflexion{-}llama{-}3.1{-}8b{-}instruct.
\vspace{24pt}

{ \footnotesize  \setcounter{utterance}{1}
\setlength{\tabcolsep}{0pt}
\begin{supertabular}{c@{$\;$}|p{.15\linewidth}@{}p{.15\linewidth}p{.15\linewidth}p{.15\linewidth}p{.15\linewidth}p{.15\linewidth}}
    \# & $\;$A & \multicolumn{4}{c}{Game Master} & $\;\:$B\\
    \hline

    \theutterance \stepcounter{utterance}  
    & & & \multicolumn{2}{p{0.3\linewidth}}{
        \cellcolor[rgb]{0.9,0.9,0.9}{
            \makecell[{{p{\linewidth}}}]{
                \texttt{\tiny{[GM$|$GM]}}
                {'hidden': {'team': ['dab', 'jet', 'pump'], 'innocent': ['hatch', 'steer', 'darn', 'counter'], 'opponent': [], 'assassin': []}, 'revealed': {'team': {'team': [], 'innocent': [], 'opponent': [], 'assassin': []}, 'opponent': {'team': [], 'innocent': [], 'opponent': [], 'assassin': []}}}
            }
        }
    }
    & & \\ \\

    \theutterance \stepcounter{utterance}  
    & & \multicolumn{4}{p{0.6\linewidth}}{
        \cellcolor[rgb]{0.9,0.9,0.9}{
            \makecell[{{p{\linewidth}}}]{
                \texttt{\tiny{[P1$\langle$GM]}}
                \texttt{You have attempted this same task before and did not succeed. These are the plans you wrote after those attempts.} \\
\texttt{Trial \#1: New plan: I got stuck because my clue was too narrow and only related to three words. What is worse, I didn't consider the potential for my partner to guess an opponent or distractor word, which would end their turn and give the opposing team a chance to make a turn. Next time, I will choose a clue word that relates to a larger subset of our team's words, and also consider the potential for my partner to guess an opponent or distractor word. I will aim to create a clue that is broad enough to be helpful, but not so broad that it becomes too vague.} \\
\\ 
\texttt{Trial \#2: New plan: In my previous attempts, I was too focused on guessing the remaining words quickly, which led me to choose clues that were too narrow and didn't consider the potential risks. What's worse, I didn't take advantage of the fact that my partner had already guessed some of our team's words, such as "cock" and "flock", which could have given me more information to work with. Next time, I will take a more strategic approach by considering the words that have already been guessed and choosing a clue that relates to a larger subset of our team's words, while also being mindful of the potential risks of my partner guessing an opponent or distractor word. I will also try to create a clue that is broad enough to be helpful, but not so broad that it becomes too vague. Additionally, I will pay closer attention to the remaining words for the opponent and distractor words, and try to create a clue that is less likely to lead my partner to guess one of those words.} \\
\\ 
\texttt{Trial \#3: New plan: In my previous attempts, I was too focused on choosing a clue that was broad enough to be helpful, but not too broad that it becomes too vague. What's worse, I didn't consider the potential for my clue to be too general and relate to the distractor words as well. Next time, I will take a more nuanced approach by choosing a clue that is specific enough to relate to a subset of our team's words, while also being mindful of the potential risks of my clue relating to the distractor words. I will also pay closer attention to the words that have already been guessed and try to create a clue that is more likely to help my partner guess our team's words, rather than the opponent's or distractor words. Additionally, I will try to create a clue that is more semantically related to the target words, rather than just being a general association. For example, if our team's words are related to water, I will choose a clue that is more directly related to water, such as "liquid", rather than a more general word like "fluid".} \\
\\ 
\\ 
\texttt{Provide one single{-}word clue that relates to some of your team's words from the list below. You can choose to how many and to which words your clue relates to.} \\
\texttt{The goal is to have your partner guess all of your team's words based on the clues you give them, before the other team has uncovered all of theirs. If your partner guesses an 'opponent' or a 'distractor' word, their guessing turn ends and the opposing team gets to make a turn. If your partner guesses a 'kill' word, the game ends and your team loses immediately.} \\
\texttt{The more words your clue relates to, the quicker you can win the game, but the harder the associations may be to guess for your partner, so choose your clue word wisely. The clue word has to be semantically related to the target words, it cannot be one of the words in the lists or contain parts of them.} \\
\\ 
\texttt{Always give your single{-}word clue and your comma{-}separated list of related target words in the following format and make your answers as short as possible, never include any other text than is required in this form:} \\
\\ 
\texttt{CLUE: <WORD>} \\
\texttt{TARGETS: <TARGETS>} \\
\\ 
\texttt{Your team words are: dab, jet, pump.} \\
\texttt{Your opponent's team words are: .} \\
\texttt{Distractor words are: hatch, steer, darn, counter.} \\
\texttt{The kill words are: .} \\
            }
        }
    }
    & \\ \\

    \theutterance \stepcounter{utterance}  
    & \multicolumn{4}{p{0.6\linewidth}}{
        \cellcolor[rgb]{0.9,0.9,0.9}{
            \makecell[{{p{\linewidth}}}]{
                \texttt{\tiny{[P1$\rangle$GM]}}
                \texttt{CLUE: fluid} \\
\texttt{TARGETS: dab, jet, pump} \\
            }
        }
    }
    & & \\ \\

    \theutterance \stepcounter{utterance}  
    & & & \multicolumn{2}{p{0.3\linewidth}}{
        \cellcolor[rgb]{0.9,0.9,0.9}{
            \makecell[{{p{\linewidth}}}]{
                \texttt{\tiny{[GM$|$GM]}}
                \texttt{fluid} \\
            }
        }
    }
    & & \\ \\

    \theutterance \stepcounter{utterance}  
    & & & \multicolumn{2}{p{0.3\linewidth}}{
        \cellcolor[rgb]{0.9,0.9,0.9}{
            \makecell[{{p{\linewidth}}}]{
                \texttt{\tiny{[GM$|$GM]}}
                ['dab', 'jet', 'pump']
            }
        }
    }
    & & \\ \\

    \theutterance \stepcounter{utterance}  
    & & & \multicolumn{2}{p{0.3\linewidth}}{
        \cellcolor[rgb]{0.9,0.9,0.9}{
            \makecell[{{p{\linewidth}}}]{
                \texttt{\tiny{[GM$|$GM]}}
                {'word': 'dab', 'assignment': 'team'}
            }
        }
    }
    & & \\ \\

    \theutterance \stepcounter{utterance}  
    & & & \multicolumn{2}{p{0.3\linewidth}}{
        \cellcolor[rgb]{0.9,0.9,0.9}{
            \makecell[{{p{\linewidth}}}]{
                \texttt{\tiny{[GM$|$GM]}}
                {'word': 'jet', 'assignment': 'team'}
            }
        }
    }
    & & \\ \\

    \theutterance \stepcounter{utterance}  
    & & & \multicolumn{2}{p{0.3\linewidth}}{
        \cellcolor[rgb]{0.9,0.9,0.9}{
            \makecell[{{p{\linewidth}}}]{
                \texttt{\tiny{[GM$|$GM]}}
                {'word': 'pump', 'assignment': 'team'}
            }
        }
    }
    & & \\ \\

    \theutterance \stepcounter{utterance}  
    & & & \multicolumn{2}{p{0.3\linewidth}}{
        \cellcolor[rgb]{0.9,0.9,0.9}{
            \makecell[{{p{\linewidth}}}]{
                \texttt{\tiny{[GM$|$GM]}}
                \texttt{opponent has won} \\
            }
        }
    }
    & & \\ \\

    \theutterance \stepcounter{utterance}  
    & \multicolumn{4}{p{0.6\linewidth}}{
        \cellcolor[rgb]{0.95,0.95,0.95}{
            \makecell[{{p{\linewidth}}}]{
                \texttt{\tiny{[P2$\rightarrow$P1]}}
                \texttt{Post{-}episode Reflexion (not sent to GM):} \\
\texttt{New plan: In my previous attempts, I was too focused on choosing a clue that was broad enough to be helpful, but not too broad that it becomes too vague. What's worse, I didn't consider the potential for my clue to be too general and relate to the distractor words as well. Next time, I will take a more nuanced approach by choosing a clue that is specific enough to relate to a subset of our team's words, while also being mindful of the potential risks of my clue relating to the distractor words. I will also pay closer attention to the words that have already been guessed and try to create a clue that is more likely to help my partner guess our team's words, rather than the opponent's or distractor words. Additionally, I will try to create a clue that is more semantically related to the target words, rather than just being a general association. For example, if our team's words are related to water, I will choose a clue that is more directly related to water, such as "liquid", rather than a more general word like "fluid".} \\
            }
        }
    }
    & & \\ \\

\end{supertabular}
}

\end{document}
